\documentclass[10pt,a4paper]{amsart}
\usepackage[margin=19mm]{geometry}
\usepackage{amsmath,amssymb,mathtools}
\usepackage[scr=rsfs,cal=euler]{mathalfa}
\usepackage{xfrac}
\usepackage[unicode,hidelinks,bookmarks=false]{hyperref}
\newcommand{\ie}{i.e.\ }
\newcommand{\cf}{cf.\ }
\newcommand{\resp}{respectively\ }
\newcommand{\Subsection}{\subsection}
\providecommand{\nobreakdash}{\nobreak\hbox{-}}
\setlength{\emergencystretch}{2em}
\multlinegap0pt
\usepackage{bm}
\usepackage{bbm}
\usepackage{dsfont}
\usepackage{xspace}
\usepackage{cancel}
\usepackage{mathtools}
\usepackage{cleveref}
\usepackage{amsthm}
\makeatletter
\@ifundefined{theorem}{\newtheorem{theorem}{Theorem}}{}
\@ifundefined{lemma}{\newtheorem{lemma}[theorem]{Lemma}}{}
\makeatother
\DeclareMathOperator*{\Argmin}{arg\,min}
\DeclareMathOperator{\Span}{span}
\newcommand{\ewc}{e^{\mathrm{wor-app}}}
\newcommand{\modone}{\bmod{}1}
\renewcommand{\hat}{\widehat} 
\title[Minimal Subsampled Rank-1 Lattices for Multivariate Approxim]{Minimal Subsampled Rank-1 Lattices for Multivariate Approximation with Optimal Convergence Rate}
\author{Felix Bartel, Alexander D. Gilbert, Frances Y. Kuo, Ian H. Sloan}
\date{Source: arXiv:2506.07729v3 (CC BY 4.0)}
\begin{document}
\maketitle

\noindent\emph{Source and licence.} Minimal Subsampled Rank-1 Lattices for Multivariate Approximation with Optimal Convergence Rate. Version arXiv:2506.07729v3 (CC BY 4.0), distributed under the Creative Commons Attribution 4.0 International licence, as recorded in the arXiv licence metadata for this exact version and shown on the abstract page https://arxiv.org/abs/2506.07729v3. Full rights notice: licenses/arxiv-2506.07729-CC-BY-4.0.txt.\par\smallskip

\noindent\textit{Editorial scope.} Counted unit: the Theorem whose source label is \texttt{withPhi} in the article (source0.tex lines 2601--2621, source label \texttt{withPhi}), delivered together with the article's own complete original proof of it (source0.tex lines 2623--2676, one \texttt{proof} environment of 54 lines). No mathematics is written, repaired or re-derived: statement and proof are the article's own text line for line. Required context retained uncounted: eq:H (lines 2208--2214); eq:lsqr (lines 2400--2403); eq:L (lines 2415--2419); generalbound (lines 2474--2490); eq:Phi (lines 2476--2481); constructions (lines 2532--2548). Every definition, display and label of those lines is retained unchanged. Omitted from this package and not redistributed: 1--2207, 2215--2399, 2404--2414, 2420--2473, 2482--2531, 2549--2600, 2622--2622, 2677--3644. Nothing inside a retained excerpt is omitted or altered: every retained body is delivered exactly as the article's lines are written, so no omission receipt of any kind accompanies this package. Citations: the retained excerpts carry no citation command at all, so no bibliography entry of the article is transcribed and no reference is redistributed. Reference adaptation: none was needed and none was made; every reference inside the delivered unit resolves inside it, so no unresolved reference or citation is produced. Printed slips kept literal: no printed word, symbol, hypothesis or number of the counted statement or of its proof was altered. Layout and class: the article's own class is \texttt{amsart} with packages including fontenc, inputenc, babel, hyperref, enumitem, graphicx, array, todonotes, xcolor, biblatex, xpatch, mathtools, amssymb, amsthm; the wrapper is the lane's pinned \texttt{amsart} editorial wrapper at 10pt on A4, which restates the theorem-like environments the retained text needs and adds the article's own macro definitions that the retained text needs (\texttt{\textbackslash Argmin}, \texttt{\textbackslash Span}, \texttt{\textbackslash ewc}, \texttt{\textbackslash hat}, \texttt{\textbackslash modone}), with extra wrapper packages (bm, bbm, dsfont, xspace, cancel, mathtools, cleveref). The delivered numbering is the unit's own. Assets: no retained span contains a figure environment, a graphics-inclusion command, a PDF-page inclusion, a table or a TikZ picture; no image file is copied into this package, the package declares no asset dependency and no image file is redistributed with it. Licence: version arXiv:2506.07729v3 of this article is distributed under the Creative Commons Attribution 4.0 International licence, as recorded in the arXiv licence metadata for this exact version and shown on the abstract page https://arxiv.org/abs/2506.07729v3. A full rights notice is delivered at licenses/arxiv-2506.07729-CC-BY-4.0.txt. Characters: every retained excerpt is pure ASCII, so the delivered engine, XeLaTeX with its default Latin Modern text font, reports no missing character.
\begin{equation}\label{eq:H}
    f\in H
    = \Big\{ f\colon\mathds T^d \to\mathds C : 
    \|f\|_H^2
    = \sum_{\bm h\in\mathds Z^d} r(\bm h) \, |\hat f_{\bm h}|^2
    < \infty \Big\} \,,
\end{equation}

\begin{equation}\label{eq:lsqr}
    S_{\mathcal B}^{\bm X_J} f
    \coloneqq \Argmin_{g\in \Span\{\exp(2\pi\mathrm i\langle\bm h,\cdot\rangle)\}_{\bm h\in\mathcal B}} \sum_{k\in J} \Big| g(\tfrac kn\bm z\modone) - f(\tfrac kn\bm z\modone) \Big|^2 \,.
\end{equation}

    \begin{equation}\label{eq:L}
        \bm L_{J,\mathcal B}
        \coloneqq \Big[ \exp(2\pi\mathrm i\tfrac kn\langle\bm h, \bm z\rangle) \Big]_{k\in J,\, \bm h\in\mathcal B}
        \in\mathds C^{|J|\times|\mathcal B|}
    \end{equation}

\begin{theorem}\label{generalbound} Let $\mathcal B\subseteq\mathds Z^d$ be a nonempty frequency index set for which the lattice $\bm X = \{\frac kn\bm z\modone\}_{k=0}^{n-1} \subseteq\mathds T^d$ has the reconstructing property.
    For $J\subseteq\{0,\dots,n-1\}$ a multiset with $|J|\ge|\mathcal B|$, suppose $\bm L_{J,\mathcal B}$ in \eqref{eq:L} has full rank, and define
    \begin{equation}\label{eq:Phi}
        \bm\Phi_{J,\mathcal B}
        \coloneqq \Big[
        \frac{\exp(2\pi\mathrm i\frac{k}{n}\langle\bm h,\bm z\rangle)}{\sqrt{r(\bm h)}}
        \Big]_{k\in J,\,\bm h\notin\mathcal B} \,.
    \end{equation}
    Then the worst-case error of the least squares approximation defined in \eqref{eq:lsqr} with respect to the function space defined in \eqref{eq:H} satisfies the bound
    \begin{equation*}
        \ewc(S_{\mathcal B}^{\bm X_J})^2
        \le \sup_{\bm h\notin\mathcal B} \frac{1}{r(\bm h)}
        + 
        \frac{ \|\bm\Phi_{J,\mathcal B}\|_{2}^{2} }{ \sigma_{\min}^2(\bm L_{J,\mathcal B}) } \,,
    \end{equation*}
    where $\sigma_{\min}(\bm L_{J, \mathcal B})$ denotes the smallest singular value of $\bm L_{J,\mathcal B}$.
\end{theorem} 

\begin{lemma}\label{constructions} Let $\mathcal B\subseteq\mathds Z^d$ be a nonempty frequency index set for which the lattice $\bm X = \{\frac kn\bm z\modone\}_{k=0}^{n-1} \subseteq\mathds T^d$ has the reconstructing property.
    Further, let $t>0$ and $J\subseteq\{0, \dots, n-1\}$ be a multiset of uniformly i.i.d.\ drawn integers with
    \begin{equation}\label{eq:12}
        |J| \ge 12 \, |\mathcal B| \, (\log|\mathcal B| + t) \,.
    \end{equation}
    For $\bm L_{J,\mathcal B}$ defined as in \eqref{eq:L} we have the following two inequalities, each holding with probability exceeding $1-\exp(-t)$,
    \begin{equation}\label{eq:sepultura}
        \frac{|J|}{2} \le \sigma_{\min}^2(\bm L_{J,\mathcal B})
        \quad\text{and}\quad
        \sigma_{\max}^2(\bm L_{J,\mathcal B}) \le \frac{3|J|}{2} \,.
    \end{equation}
    For $|J|\ge 3$, $\bm\Phi_{J,\mathcal B}$ as in \eqref{eq:Phi}, and $\bm\Phi_{\mathcal B} = \bm\Phi_{\{0,\dots,n-1\},\mathcal B}$, we have with probability exceeding $1-2^{3/4}\exp(-t)$ that
    \begin{equation}\label{eq:watain}
        \|\bm\Phi_{J,\mathcal B}\|_{2}^2
        \le 21 \, (\log|J|+t) \sum_{\bm h\notin\mathcal B} \frac{1}{r(\bm h)} + 2\frac{|J|}{n} \|\bm\Phi_{\mathcal B}\|_{2}^2 \,.
    \end{equation}
\end{lemma} 

\begin{theorem}\label{withPhi} Let $\mathcal B\subseteq\mathds Z^d$ be a nonempty frequency index set for which the lattice $\bm X = \{\frac kn\bm z\modone\}_{k=0}^{n-1} \subseteq\mathds T^d$ has the reconstructing property.
    Further, let $t\ge4$ and $J\subseteq\{0, \dots, n-1\}$ be a multiset of uniformly i.i.d.\ drawn integers with
    \begin{equation*}
|J| := \lceil 12 \, |\mathcal B| \, (\log|\mathcal B| + t)\rceil\,.
    \end{equation*}
    Then the least squares approximation defined in \eqref{eq:lsqr} using samples on the subsampled lattice $\bm X_J$ satisfies with probability exceeding $1-3\exp(-t)$ that
    \begin{equation*}
        \ewc(S_{\mathcal B}^{\bm X_J})^2
        \le \sup_{\bm h\notin\mathcal B} \frac{1}{r(\bm h)}
        + 
        \frac{42\,(\log|J|+t)}{|J|}
\sum_{\bm h\notin\mathcal B} \frac{1}{r(\bm h)} + \frac{4}{n} \|
        \bm\Phi_{\mathcal B}\|_{2}^2 \,,
    \end{equation*}
    with respect to the function space defined in \eqref{eq:H} and $\bm\Phi_{\mathcal B} = \bm\Phi_{\{0,\dots,n-1\},\mathcal B}$ as in \eqref{eq:Phi}. Additionally we have
    \begin{equation}\label{eq:JB}
         \frac{6\,(\log|J|+t)}{|J|} 
         \le \frac{1}{|\mathcal B|} 
         \le \frac{13\,(\log|J|+t)}{|J|} \,.
    \end{equation}
\end{theorem} 

\begin{proof} We use the bound of the smallest singular value of $\bm L_{J,\mathcal B}$ and the bound on the spectral norm of $\bm\Phi_{J,\mathcal B}$ of \Cref{constructions}.
    By Boole's inequality, the overall probability exceeds $1-\exp(-t)-2^{3/4}\exp(-t) \ge 1-3\exp(-t)$, which is the claimed probability.
    Plugging in the spectral bounds into \Cref{generalbound}, we obtain
    \begin{equation*}
        \ewc(S_{\mathcal B}^{\bm X_J})^2
        \le \sup_{\bm h\notin\mathcal B} \frac{1}{r(\bm h)}
        + 
        \frac{2}{|J|} \Big( 21 \,(\log|J|+t) \sum_{\bm h\notin\mathcal B} \frac{1}{r(\bm h)} + 2\frac{|J|}{n} \|\bm\Phi_{\mathcal B}\|_{2}^2\Big) \,.
    \end{equation*}
    It remains to estimate the scaling factor for the sum.

    We first show that for $b\ge 1$ and $t\ge 4$ we have
    \begin{equation}\label{eq:hard}
        12\,b\,(\log b+t)
        \le b^2\exp(t) \,.
    \end{equation}
    Since $\log b \le b-1$ this follows from the stronger assertion
    \begin{equation*}
        12\,(b-1+t)
        \le b\exp(t)
        \quad\Leftrightarrow\quad
        (\exp(t)-12)\,b +12-12t
        \ge 0 \,.
    \end{equation*}
    The left-hand side is non-decreasing in $b$ since $t\ge 4 \ge \log(12)$.
    With $b\ge 1$, it remains to check the non-negativity for $b=1$, which is satisfied for $t\ge 4$ showing~\eqref{eq:hard}.

    We are ready to prove an upper bound on $(\log j+t)/j$ in terms of $1/b$ when $j := \lceil\,12\,b\,(\log b+t)\,\rceil \ge 12\,b\,(\log b+t)$.
    Since $(\log j+t)/j$ is increasing with decreasing $j$, it suffices to show an estimate for $\ell\le j$ such that
    \begin{equation*}
        12\,b\,(\log b+t)
        \le \ell
        \le b^2\exp(t) \,.
    \end{equation*}
    Such a value of $\ell\in\mathds R$ exists due to \eqref{eq:hard}.
    Then it follows that
    \begin{equation*}
        \frac{\log j+t}{j}
        \le \frac{\log\ell+t}{\ell}
        \le \frac{2 \, (\log b + t)}{12 \, b \, (\log b + t)}
        = \frac{1}{6\,b} \,.
    \end{equation*}
    
    Finally we prove an upper bound on $1/b$ in terms of $(\log j+t)/j$.
    Since $b\ge 1$ and $t\ge 4$, we have $12 \, b \, (\log b+t) \ge \frac{48}{49} \lceil 12 \, b \, (\log b+t)\rceil$.
    Thus
    \begin{equation*}
        \frac{1}{b}
        \le \frac{49}{48} \frac{12 \, (\log b+t)}{\lceil 12 \, b \, (\log b+t)\rceil}
        \le \frac{13 \, (\log j+t)}{j} \,,
    \end{equation*}
    where we used $j > 12\,b \ge b$ for $\log b < \log j$ and $49\cdot 12/48 <13$.
    This completes the proof.
\end{proof} 

\end{document}
